\documentclass[11pt,a4paper]{article}

\usepackage[T1]{fontenc}
\usepackage[utf8]{inputenc}
\usepackage{lmodern}
\usepackage{microtype}
\usepackage[margin=1in]{geometry}
\usepackage{xcolor}
\usepackage{listings}
\usepackage{booktabs}
\usepackage{array}
\usepackage{tabularx}
\usepackage{enumitem}
\usepackage{fancyhdr}
\usepackage[hidelinks]{hyperref}

% ---- Colours and code style ----
\definecolor{codebg}{gray}{0.96}
\definecolor{codeframe}{gray}{0.80}

\lstdefinestyle{davinci}{
  basicstyle=\ttfamily\small,
  backgroundcolor=\color{codebg},
  frame=single,
  rulecolor=\color{codeframe},
  breaklines=true,
  breakatwhitespace=false,
  columns=fullflexible,
  keepspaces=true,
  showstringspaces=false,
  xleftmargin=6pt,
  xrightmargin=6pt,
  aboveskip=8pt,
  belowskip=8pt
}
\lstset{style=davinci}

% ---- Header and footer ----
\pagestyle{fancy}
\fancyhf{}
\fancyhead[L]{\footnotesize DAVINCI}
\fancyhead[R]{\footnotesize D4: Implementation Report}
\fancyfoot[C]{\thepage}
\renewcommand{\headrulewidth}{0.4pt}

\setlength{\parskip}{0.6em}
\setlength{\parindent}{0pt}

\begin{document}

% ================= Title block (matches D2/D3) =================
\thispagestyle{plain}
\begin{center}
  {\LARGE\bfseries D4: Implementation Report of a Secure\\[3pt]
   Deep Attestation Framework}\\[12pt]
  {\large DAVINCI: Deep Attestation Framework for Virtualized\\[2pt]
   Network Communication Integrity}\\[18pt]
\end{center}

\begin{center}
\begin{tabular}{c}
University of Padua Technical Leader: Prof. Mauro Conti\\[3pt]
Orange Technical Leader: Dr. Ghada Arfaoui\\[3pt]
Authors: Rahul Saha, Gulshan Kumar, Marco Pasca\\[3pt]
Reviewers: Ghada Arfaoui and Thibaut Jacques, Orange\\
\end{tabular}
\end{center}

\vspace{12pt}

\begin{center}
\begin{tabularx}{\textwidth}{l l X l}
\toprule
\textbf{Version} & \textbf{Date} & \textbf{Authors} & \textbf{Status}\\
\midrule
v1.0 & 10/07/2026 & Rahul Saha, Gulshan Kumar & Submitted\\
\bottomrule
\end{tabularx}
\end{center}

\vspace{10pt}
\hrule
\vspace{4pt}

% ================= 1. Executive Summary =================
\section{Executive Summary}

This deliverable explains how to set up, build, and run the DAVINCI prototype,
and how to read the results it produces. It follows the design given in D3 and
implements the attestation scheme and the security properties defined in D2.
Where D3 describes what each component should do, this document gives the
concrete commands that get the framework running on a real machine, together
with a worked example that exercises the complete workflow from evidence
collection to the final trust decision.

The document has four parts. The first part prepares the environment where
DAVINCI runs, either inside a virtual machine created with Vagrant and Ansible,
or directly on a macOS or Linux host. The second part sets up DAVINCI itself,
meaning the key material and the build of the three services. The third part
runs a complete example. The fourth part explains the output of that example and
states the current limitations of the prototype.

The prototype consists of three services written in Rust and a small set of
Python scripts that act as evidence sources for the demonstration. The Privacy
Control Engine on port 3032 applies the privacy transformations to evidence. The
Attestation Manager on port 3030 checks each piece of evidence, builds the
measurement chain and the trust graph, and exports a signed result. The Verifier
on port 3031 checks the signed result and produces the trust decision. All
signatures use ECDSA over the NIST P256 curve and all hashing uses SHA-256,
which are the cryptographic primitives assumed in D2.

% ================= 2. System Overview =================
\section{System Overview}

The concrete DAVINCI instantiation maps the five abstract algorithms of the
scheme onto the three services as follows. Setup is the generation of the
Attestation Manager key pair, whose public key becomes the trust anchor held by
the Verifier. Collect is the per source evidence gathering, where each source
signs its measurement with a Root of Trust key and the Attestation Manager
verifies that signature before admitting the evidence. Transform is the privacy
processing performed by the Privacy Control Engine. Attest is the construction of
the hash linked measurement chain and the signed trust graph. Verify is the
Verifier check that combines the signature verification, the freshness check, and
the graph validity check into a single decision.

A single attestation run therefore flows through the services in one direction.
Evidence enters the Attestation Manager, is authenticated and privacy processed,
becomes a node in the measurement chain and a vertex in the trust graph with
signed edges, and is finally exported as a signed output that the Verifier
appraises against its trust anchor.

% ================= 3. Environment =================
\section{Setting up the Environment}

There are two supported ways to prepare a machine. The first uses a virtual
machine and is the most reproducible, because every dependency is installed the
same way every time. The second installs the same dependencies directly on macOS
or Linux, which is convenient when you already work on the machine day to day.

\subsection{A virtual machine with Vagrant and Ansible}

This option builds a clean Ubuntu 22.04 virtual machine and provisions it
automatically. You need VirtualBox and Vagrant installed on your host. On macOS
you can install both with Homebrew.

\begin{lstlisting}
brew install --cask virtualbox
brew install --cask vagrant
\end{lstlisting}

On a Debian or Ubuntu host you can install them with apt.

\begin{lstlisting}
sudo apt update
sudo apt install -y virtualbox vagrant
\end{lstlisting}

The repository already contains a \texttt{Vagrantfile} and an Ansible playbook at
\texttt{provisioning/}\allowbreak\texttt{playbook.yml}. From the root of the DAVINCI directory,
create and provision the machine with a single command.

\begin{lstlisting}
vagrant up
\end{lstlisting}

This downloads the Ubuntu image, boots the machine, and then runs the Ansible
playbook inside it. The playbook installs the build tools and the OpenSSL
development libraries, installs the Rust toolchain, installs Python, generates
the DAVINCI keys, and builds the three services. The three service ports 3030,
3031, and 3032 are forwarded to the same ports on your host, so you can reach the
services from your own terminal. When provisioning finishes, log in to the
machine and move to the project directory, which is mounted inside the machine at
\texttt{/vagrant}.

\begin{lstlisting}
vagrant ssh
cd /vagrant
\end{lstlisting}

You are now ready for the example in Section 5. To provision again without
rebuilding the machine, run \texttt{vagrant provision}. To shut the machine down,
run \texttt{vagrant halt}, and to delete it completely, run
\texttt{vagrant destroy}. If you prefer to run the same playbook against a plain
Ubuntu host instead of a virtual machine, install Ansible and run it in local
mode from the project directory.

\begin{lstlisting}
sudo apt install -y ansible
ansible-playbook -i "localhost," -c local provisioning/playbook.yml
\end{lstlisting}

\subsection{Setting up directly on macOS or Linux}

If you would rather run everything on your own machine, install four things: the
Rust toolchain, OpenSSL, pkg-config, and Python 3. Install the Rust toolchain
through rustup, which does not require administrator rights.

\begin{lstlisting}
curl --proto '=https' --tlsv1.2 -sSf https://sh.rustup.rs | sh -s -- -y
source "$HOME/.cargo/env"
\end{lstlisting}

On macOS, install OpenSSL and pkg-config with Homebrew, then point pkg-config at
the OpenSSL you just installed so that the Rust build can find it.

\begin{lstlisting}
brew install openssl@3 pkg-config
export PKG_CONFIG_PATH="$(brew --prefix openssl@3)/lib/pkgconfig"
\end{lstlisting}

On Ubuntu or Debian, install the equivalent packages with apt.

\begin{lstlisting}
sudo apt update
sudo apt install -y build-essential pkg-config libssl-dev openssl python3
\end{lstlisting}

Python 3 is already present on macOS and on most Linux systems. The demonstration
scripts use only the Python standard library and the openssl command line tool,
so there are no extra Python packages to install. You can confirm the toolchain
is ready with the following commands.

\begin{lstlisting}
rustc --version
openssl version
python3 --version
\end{lstlisting}

% ================= 4. DAVINCI setup =================
\section{Setting up DAVINCI}

With the environment in place, move to the root of the DAVINCI directory. Inside
the virtual machine this is \texttt{/vagrant}. The \texttt{attestation-manager},
\texttt{verifier}, and \texttt{pce-service} folders each hold one Rust service.
The \texttt{keys} folder holds the key material. The \texttt{demo} folder holds
the example scripts. The \texttt{docs} folder holds the design and implementation
reports.

\subsection{Generating the keys}

DAVINCI relies on two kinds of keys. The Attestation Manager has its own key
pair. It signs its results with the private key, and the Verifier trusts the
matching public key, which plays the role of the trust anchor in D2 and D3. There
is also one key pair for each Root of Trust source, such as the TPM and the
confidential computing report. In a real deployment these would be hardware
protected keys inside the attestation hardware. In the prototype they are
software keys that stand in for that hardware, so that the signature checks are
genuine even without special hardware. Generate all of the keys with the helper
script.

\begin{lstlisting}
./keys/gen-keys.sh
\end{lstlisting}

This writes the Attestation Manager key pair and the per source key pairs into
the keys folder. The private keys are never committed to version control, so this
step must be run once on any fresh copy of the repository.

\subsection{Building the services}

The simplest path is the Makefile target that generates the keys and then builds
all three services.

\begin{lstlisting}
make build
\end{lstlisting}

On macOS, remember to export \texttt{PKG\_CONFIG\_PATH} first, as shown in
Section 3, so that the OpenSSL library is found during the build. The first build
downloads and compiles the dependencies and can take a few minutes. Later builds
are fast. If you prefer to build the services one at a time, you can build each
crate on its own.

\begin{lstlisting}
cd pce-service && cargo build && cd ..
cd attestation-manager && cargo build && cd ..
cd verifier && cargo build && cd ..
\end{lstlisting}

\subsection{Configuration}

The services read a few settings from environment variables. The defaults are
chosen so that everything works on one machine without any configuration, so you
can usually ignore this section for a first run.

\begin{center}
\footnotesize
\begin{tabularx}{\textwidth}{>{\raggedright\arraybackslash}p{0.30\textwidth} >{\raggedright\arraybackslash}p{0.24\textwidth} X}
\toprule
\textbf{Variable} & \textbf{Default} & \textbf{Purpose}\\
\midrule
\texttt{DAVINCI\_KEYS\_DIR} & \texttt{../keys} & Folder that holds the key material.\\
\texttt{PCE\_URL} & \texttt{http://localhost:3032} & Where the Attestation Manager reaches the Privacy Control Engine.\\
\texttt{ATTESTATION\_MANAGER\_URL} & \texttt{http://localhost:3030} & Where the Verifier reaches the Attestation Manager.\\
\texttt{DATABASE\_URL} & unset & PostgreSQL connection for the Privacy Control Engine. If unset, it keeps data in memory.\\
\texttt{AM\_BIND} & \texttt{0.0.0.0:3030} & Listen address of the Attestation Manager.\\
\texttt{VERIFIER\_BIND} & \texttt{0.0.0.0:3031} & Listen address of the Verifier.\\
\texttt{PCE\_BIND} & \texttt{0.0.0.0:3032} & Listen address of the Privacy Control Engine.\\
\bottomrule
\end{tabularx}
\end{center}

% ================= 5. Running an example =================
\section{Running an Example Instance}

The example drives the whole workflow. It sends platform, hypervisor, guest, and
runtime evidence to the Attestation Manager, asks the Verifier for a decision,
and prints the result. It runs two scenarios one after the other. The first uses
authentic evidence and should be trusted. The second tampers with the guest
evidence after the Root of Trust has signed it, and should be rejected.

\subsection{Starting the services}

Start the three services in the correct order. Start the Privacy Control Engine
first, because the Attestation Manager depends on it. Give each service its own
terminal, because each one runs in the foreground and prints a log.

\begin{lstlisting}
make run-pce                  # terminal 1, port 3032
make run-attestation-manager  # terminal 2, port 3030
make run-verifier             # terminal 3, port 3031
\end{lstlisting}

You can confirm that all three are up with three quick requests. The first
returns a status of ok, the second returns the Attestation Manager public key in
PEM form, and the third returns a status of ok.

\begin{lstlisting}
curl -s http://127.0.0.1:3032/health
curl -s http://127.0.0.1:3030/pubkey
curl -s http://127.0.0.1:3031/health
\end{lstlisting}

\subsection{Running the example}

In a fourth terminal, move to the demo folder and run the comparison example,
which runs the trusted scenario and then the tampered scenario.

\begin{lstlisting}
cd demo
python3 demo_comparison.py
\end{lstlisting}

There is also a single end to end run that only performs the trusted scenario,
available as \texttt{python3 demo\_scenario.py}. The same runs are available
through the Makefile as \texttt{make demo-compare} and \texttt{make demo}.

% ================= 6. Output and limitations =================
\section{Reading the Output and Current Limitations}

\subsection{The trusted scenario}

In the trusted scenario every layer sends authentic evidence, so the Attestation
Manager admits each item, builds the graph, and signs it. The Verifier then
confirms every check and returns a trusted decision.

\begin{lstlisting}
Trust Decision: TRUSTED
Checks:
  PASS Output Signature (ECDSA-P256)
  PASS Freshness (challenge binding)
  PASS Platform Root Present
  PASS Edge Signatures (ECDSA-P256)
  PASS Layer Monotonicity
  PASS Binding Policy
  PASS Acyclicity
  PASS Cross-Layer Trust Continuity
\end{lstlisting}

\subsection{The tampered scenario}

In the tampered scenario the guest measurement is changed after the Root of Trust
signed it. The Attestation Manager checks the signature over the guest
measurement, finds that it no longer matches, and rejects that evidence. The
guest node therefore never enters the graph. The runtime node above it then has
no path down to the platform, so the Verifier reports that the cross layer trust
continuity check has failed and returns a not trusted decision.

\begin{lstlisting}
3. Guest SEV-SNP evidence ...  (TAMPERED after RoT signature)
   AM: rejected_auth: RoT signature invalid, evidence not authentic (Auth = 0)

Trust Decision: NOT TRUSTED
  PASS Output Signature (ECDSA-P256)
  PASS Freshness (challenge binding)
  PASS Platform Root Present
  PASS Edge Signatures (ECDSA-P256)
  PASS Layer Monotonicity
  PASS Binding Policy
  PASS Acyclicity
  FAIL Cross-Layer Trust Continuity  (Application node has no path to Platform)
\end{lstlisting}

The important point is that this decision comes from real signature checks and
graph checks, and not from any label attached to the session.

\subsection{The signed output and the trust graph}

The Attestation Manager exports the result at the endpoint
\texttt{GET /attest/output/\{session\}}, which takes the challenge as a
\texttt{nonce} query parameter. The response contains the output and a proof,
which is the Attestation Manager signature over that output. A trimmed example
follows.

\begin{lstlisting}
{
  "out": {
    "sid": "d4_example_1783669785",
    "challenge": "chal_demo_0001",
    "ts_out": 1783669785,
    "manifest": { "redacted": 8, "aggregated": 0, "retained": 1 }
  },
  "pf": "304602210082649319d58a1a7ae4d97faa074fe3cefd343c..."
}
\end{lstlisting}

The manifest records what the Privacy Control Engine did to the evidence. In this
run it removed eight sensitive fields across the four evidence items, which are
the host name and the address field of each item, and kept the fields that
verification needs. Each node in the graph carries its profile, its measurement,
its position in the hash chain, and its hash.

\begin{lstlisting}
{
  "jti": "645f3ed0-e43c-4334-a3e0-4f8c0182dd5d",
  "eat_profile": "layer:Application,mode:Runtime,source:Ima",
  "eat_meas": "{ \"claims\": {}, \"measurement\": \"bf1f8f7c3bba...\" }",
  "iat": 1783669785,
  "x_davinci_prev": "f792985febed579dd70c0e110474e59268dedc84717a19255fa9ea9496e93859",
  "hash": "f9a9e69f02e52dcbf3ba7c54b5c341983aa0c3fadbe23e64ae2b3bebf5e980ef"
}
\end{lstlisting}

Each edge records a trust dependency and carries the Attestation Manager
signature over that edge.

\begin{lstlisting}
{
  "from": "4165ca62-aa7b-46c8-ae29-5fe7566a7e65",
  "to": "5bec0611-88bf-4992-8ca7-a3e492f3c2c2",
  "relationship": "VmIncludesHypervisor",
  "signature": "30440220429844689d9cf0dfda5a918ebb3592e794a1c929..."
}
\end{lstlisting}

\subsection{How the checks map to the security properties}

Each check the Verifier performs corresponds to a property defined in D2.

\begin{center}
\begin{tabularx}{\textwidth}{>{\raggedright\arraybackslash}p{0.38\textwidth} X}
\toprule
\textbf{Verifier check} & \textbf{D2 property}\\
\midrule
Output Signature, Edge Signatures & Authenticity, since only the Attestation Manager private key can produce them.\\
Freshness & Freshness and replay resistance, through the fresh challenge bound into the signed output.\\
Platform Root, Layer Monotonicity, Acyclicity, Binding Policy, Cross-Layer Trust Continuity & Deep integrity and cross layer trust continuity, through a signed and well formed path from every layer to the platform anchor.\\
Rejection of tampered evidence at the Attestation Manager & Evidence authenticity that underlies multi modal soundness.\\
Privacy Control Engine transform & The privacy properties, through pseudonymisation and field suppression.\\
\bottomrule
\end{tabularx}
\end{center}

The binding policy check corresponds to the admissible edge rule $\Pi_{\mathrm{bind}}$
of D3, which permits only the expected parent relationships between layers.

\subsection{The application interfaces}

The Attestation Manager accepts boot evidence at \texttt{POST /attest/boot} and
runtime evidence at \texttt{POST /attest/runtime}. It returns the signed output at
\texttt{GET /attest/output/\{session\}}, the raw graph for inspection at
\texttt{GET /arg/\{session\}}, and its public key at \texttt{GET /pubkey}. The
Verifier accepts a request at \texttt{POST /verify} with a session and a challenge
and returns the full decision with the per check details. The Privacy Control
Engine accepts \texttt{POST /privacy/transform},
\texttt{POST /validate/node}, and \texttt{POST /resolve/pseudonym}, and reports
its state at \texttt{GET /health}.

\subsection{Current limitations}

The prototype is faithful to the design in the parts it implements. The following
points describe what it does not yet do, so that results are read correctly.

First, the hardware Roots of Trust are simulated. The evidence is signed with
software ECDSA keys that stand in for hardware protected keys, so the signature
checks are real, but no physical TPM or confidential computing hardware is
involved.

Second, the bootstrap attestation of the Attestation Manager host, described in
D3, is not implemented. The Verifier trusts the Attestation Manager public key
from its key folder rather than first appraising the host that holds it.

Third, the Verifier validates the graph it receives, but it does not yet require
a fixed set of layers to be present. A result that contains a smaller but
internally valid graph is accepted as far as it goes. Enforcing a required set of
layers per session is the natural next step and follows directly from the session
policy described in D3.

\subsection{Stopping the services}

Each service runs in the foreground, so stopping it is a matter of pressing Ctrl
and C in its terminal. If you started the services in the background, you can stop
them by name.

\begin{lstlisting}
pkill -f target/debug/pce-service
pkill -f target/debug/attestation-manager
pkill -f target/debug/verifier
\end{lstlisting}

% ================= 7. Conclusion =================
\section{Conclusion}

This report has shown how to prepare an environment for DAVINCI, how to build the
three services and their key material, and how to run a complete example that
produces a cryptographically justified trust decision. The worked example
demonstrates that authentic evidence is accepted and that tampered evidence is
detected and rejected, in both cases through real signature and graph checks
rather than any external label. Together with the design in D3 and the theory in
D2, this completes the description of the DAVINCI prototype and provides a
reproducible basis for further development.

\end{document}
